\documentclass{article}
\usepackage{amsmath,amssymb,amsthm}
\usepackage{algorithm}
\usepackage[noend]{algpseudocode}
\usepackage{fullpage}
\newtheorem{theorem}{Theorem}
\newtheorem{lemma}{Lemma}
\newcommand{\EE}{\mathbb{E}}
\newcommand{\RR}{\mathbb{R}}
\newcommand{\norm}[1]{\left\| #1\right\|}
\newcommand{\inner}[2]{\left\langle #1,#2\right\rangle}
\begin{document}

\section*{Algorithm Name}
\textbf{AdamW} – Adam with decoupled weight decay (non‑convex, $L$‑smooth setting).

\section*{Mathematical Formulation}
Let $f:\RR^{d}\rightarrow\RR$ be the (possibly stochastic) objective.
At iteration $t\ge 1$ we have:
\[
\begin{aligned}
g_t &:= \nabla f(w_t;\xi_t) \qquad\text{(stochastic gradient on mini‑batch }\xi_t)\\[4pt]
m_t &= \beta_1 m_{t-1} + (1-\beta_1) g_t, \\
v_t &= \beta_2 v_{t-1} + (1-\beta_2) g_t\odot g_t, \qquad (\odot\text{ elementwise product})\\[4pt]
\hat m_t &= \dfrac{m_t}{1-\beta_1^{\,t}},\qquad 
\hat v_t = \dfrac{v_t}{1-\beta_2^{\,t}};\\[4pt]
\tilde w_t &= w_t - \eta\lambda w_t \qquad\text{(decoupled weight decay)}\\[4pt]
w_{t+1} &= \tilde w_t - \eta\;\frac{\hat m_t}{\sqrt{\hat v_t}+\varepsilon},
\end{aligned}
\]
where
\begin{itemize}
    \item $\eta>0$ is the learning rate,
    \item $\lambda\ge0$ is the weight‑decay coefficient,
    \item $\beta_1,\beta_2\in[0,1)$ are exponential‑moving‑average coefficients,
    \item $\varepsilon>0$ is a small constant for numerical stability.
\end{itemize}
All operations are elementwise for vectors.

\section*{Pseudocode}
\begin{algorithm}[H]
\caption{AdamW}
\begin{algorithmic}[1]
\Require Learning rate $\eta$, weight decay $\lambda$, $\beta_1,\beta_2\in[0,1)$, $\varepsilon>0$,
        initialization $w_0\in\RR^{d}$, $m_0=0$, $v_0=0$, maximum iterations $T$.
\For{$t=1$ \textbf{to} $T$}
    \State Sample mini‑batch $\xi_t$ and compute stochastic gradient $g_t=\nabla f(w_t;\xi_t)$
    \State $m_t \gets \beta_1 m_{t-1} + (1-\beta_1)g_t$
    \State $v_t \gets \beta_2 v_{t-1} + (1-\beta_2) g_t\odot g_t$
    \State $\hat m_t \gets m_t/(1-\beta_1^{t})$
    \State $\hat v_t \gets v_t/(1-\beta_2^{t})$
    \State $\tilde w_t \gets w_t - \eta\lambda w_t$ \Comment{decoupled weight decay}
    \State $w_{t+1} \gets \tilde w_t - \eta\; \hat m_t/(\sqrt{\hat v_t}+\varepsilon)$
\EndFor
\State \Return $w_{T}$
\end{algorithmic}
\end{algorithm}

\section*{Dry Run Example}
Objective $f(w)=(w-3)^2$, thus $\nabla f(w)=2(w-3)$.  
Choose scalar hyper‑parameters: $\eta=0.1$, $\lambda=0.01$, $\beta_1=0.9$, $\beta_2=0.999$, $\varepsilon=10^{-8}$.  
Initialize $w_0=0$, $m_0=0$, $v_0=0$.

\begin{center}
\begin{tabular}{c|c|c|c|c|c}
$t$ & $g_t$ & $m_t$ & $v_t$ & $w_{t}$ & $f(w_t)$\\\hline
0 & – & 0 & 0 & $0$ & $9$\\
1 & $2(0-3)=-6$ &
  $0.9\cdot0+0.1(-6)=-0.6$ &
  $0.999\cdot0+0.001(36)=0.036$ &
  $w_1 = (1-\eta\lambda)w_0 - \eta\frac{m_1}{\sqrt{v_1}+\varepsilon}
        =0 -0.1\frac{-0.6}{\sqrt{0.036}+10^{-8}}
        \approx 1.0$ &
  $f(1)=4$\\
2 & $2(1-3)=-4$ &
  $0.9(-0.6)+0.1(-4)=-0.94$ &
  $0.999(0.036)+0.001(16)=0.051964$ &
  $w_2 = (1-0.001)1.0 -0.1\frac{-0.94}{\sqrt{0.051964}} 
        \approx 1.0+0.410\approx 1.410$ &
  $f(1.410)\approx 2.52$\\
3 & $2(1.410-3)=-3.18$ &
  $0.9(-0.94)+0.1(-3.18)=-1.146$ &
  $0.999(0.051964)+0.001(10.1124)=0.062045$ &
  $w_3 = (1-0.001)1.410 -0.1\frac{-1.146}{\sqrt{0.062045}}
        \approx 1.408+0.459\approx 1.867$ &
  $f(1.867)\approx 1.28$
\end{tabular}
\end{center}
The objective value decreases from $9$ to $\approx1.28$ in three steps, illustrating the effect of the adaptive step and weight decay.

\section*{Assumptions}
\begin{enumerate}
\item[(A1)] \textbf{$L$‑smoothness.} $f$ is differentiable and for all $x,y\in\RR^{d}$
\[
f(y)\le f(x)+\inner{\nabla f(x)}{y-x}+\frac{L}{2}\norm{y-x}^2 .
\]
\item[(A2)] \textbf{Unbiased stochastic gradients.}
\[
\EE_{\xi_t}[g_t\mid w_t]=\nabla f(w_t),\qquad\forall t.
\]
\item[(A3)] \textbf{Bounded variance.}
There exists $\sigma^2<\infty$ such that
\[
\EE_{\xi_t}\!\left[\norm{g_t-\nabla f(w_t)}^2\mid w_t\right]\le\sigma^2,\qquad\forall t.
\]
\item[(A4)] \textbf{Bounded second moments.}
For a constant $G>0$,
\[
\EE_{\xi_t}\!\left[\norm{g_t}^2\mid w_t\right]\le G^2,\qquad\forall t.
\]
\item[(A5)] \textbf{Hyper‑parameter regime.}
\[
\beta_1<\sqrt{\beta_2},\qquad
\eta\le \frac{1-\beta_1}{L\,(1+\sqrt{\beta_2})}.
\]
\end{enumerate}

\section*{Main Convergence Theorem}
\begin{theorem}[Non‑convex convergence of AdamW]\label{thm:adamw}
Under (A1)–(A5) and for any $T\ge1$, the iterates produced by AdamW satisfy
\[
\frac{1}{T}\sum_{t=1}^{T}\EE\!\left[\norm{\nabla f(w_t)}^2\right]
\;\le\;
\underbrace{\frac{2\bigl(f(w_1)-f^\star\bigr)}{\eta(1-\beta_1)T}}_{\text{initial gap term}}
\;+\;
\underbrace{\frac{L\eta G^2}{(1-\beta_1)(1-\beta_2)}\frac{\log T}{\sqrt{T}}}_{\text{adaptive term}}
\;+\;
\underbrace{2\lambda^2 G^2}_{\text{weight‑decay bias}} .
\]
Consequently, choosing $\eta = \Theta\!\bigl(1/\sqrt{T}\bigr)$ yields
\[
\frac{1}{T}\sum_{t=1}^{T}\EE\!\left[\norm{\nabla f(w_t)}^2\right]=
O\!\left(\frac{\log T}{\sqrt{T}}\right).
\]
\end{theorem}

\section*{Theoretical Properties}
\begin{itemize}
\item \textbf{Stability.} The bias‑corrected moments guarantee that the effective step size
$\eta\hat m_t/(\sqrt{\hat v_t}+\varepsilon)$ remains bounded, preventing explosion even when gradients are large.
\item \textbf{Hyper‑parameter sensitivity.}
Condition (A5) shows that $\beta_1$ must be sufficiently smaller than $\sqrt{\beta_2}$; otherwise the bound on the adaptive term deteriorates.
\item \textbf{Comparison to baselines.}
When $\lambda=0$, the bound reduces to the classical Adam guarantee (see e.g. \cite{reddi2019convergence}). The additional $2\lambda^2 G^2$ term is the price of decoupled weight decay and is independent of $T$, showing that weight decay does not affect the asymptotic rate.
\end{itemize}

\section*{Discussion}
The theorem demonstrates that AdamW inherits the $O(\log T/\sqrt{T})$ convergence rate of Adam for smooth non‑convex objectives while cleanly separating the regularization effect of weight decay. The bound is deterministic up to the stochastic variance $\sigma^2$ hidden in $G^2$; tighter analysis (e.g., using $\sigma$ directly) yields the same order but with improved constants.

\section*{Preliminaries (Definitions and Lemmas)}
\begin{lemma}[Bias‑correction bounds]\label{lem:bias}
For all $t\ge1$,
\[
\frac{1}{1-\beta_1^{\,t}}\le \frac{1}{1-\beta_1},\qquad
\frac{1}{1-\beta_2^{\,t}}\le \frac{1}{1-\beta_2}.
\]
\end{lemma}
\begin{proof}
Since $0<\beta_i<1$, the sequences $\beta_i^{\,t}$ are decreasing, hence $1-\beta_i^{\,t}\ge 1-\beta_i$ and the claim follows. \qedhere
\end{proof}

\begin{lemma}[Moment growth]\label{lem:moments}
Under (A4) and (A5),
\[
\EE\!\left[\norm{\hat m_t}^2\right]\le \frac{G^2}{(1-\beta_1)^2},
\qquad
\EE\!\left[\norm{\hat v_t}\right]\le \frac{G^2}{1-\beta_2}.
\]
\end{lemma}
\begin{proof}
From the recursion $m_t=\beta_1 m_{t-1}+(1-\beta_1)g_t$ and using (A4),
\[
\EE[\norm{m_t}^2]\le \beta_1^2\EE[\norm{m_{t-1}}^2]+(1-\beta_1)^2G^2.
\]
Unfolding the recurrence yields $\EE[\norm{m_t}^2]\le \frac{(1-\beta_1)G^2}{1-\beta_1}=G^2/(1-\beta_1)$. Applying bias correction (Lemma~\ref{lem:bias}) gives the bound for $\hat m_t$. The bound for $\hat v_t$ follows analogously using the elementwise recursion for $v_t$. \qedhere
\end{proof}

\section*{Main Proof (Step‑by‑Step Derivation)}
We prove Theorem~\ref{thm:adamw}. Throughout the proof expectations are taken w.r.t. the randomness up to iteration $t$ unless otherwise noted.

\subsection*{[Step 1] Smoothness inequality}
From (A1) with $x=w_t$, $y=w_{t+1}$,
\[
f(w_{t+1})\le f(w_t) + \inner{\nabla f(w_t)}{w_{t+1}-w_t}
+ \frac{L}{2}\norm{w_{t+1}-w_t}^2 . \tag{1}
\]

\subsection*{[Step 2] Expand the update}
Recall $w_{t+1}=w_t-\eta\lambda w_t-\eta\frac{\hat m_t}{\sqrt{\hat v_t}+\varepsilon}$.
Define
\[
\Delta_t:= w_{t+1}-w_t = -\eta\lambda w_t -\eta\frac{\hat m_t}{\sqrt{\hat v_t}+\varepsilon}. \tag{2}
\]

\subsection*{[Step 3] Plug $\Delta_t$ into (1)}
\[
\begin{aligned}
f(w_{t+1}) &\le f(w_t)
+ \inner{\nabla f(w_t)}{\Delta_t}
+ \frac{L}{2}\norm{\Delta_t}^2 .
\end{aligned} \tag{3}
\]

\subsection*{[Step 4] Separate weight‑decay and adaptive parts}
Using linearity of inner product,
\[
\inner{\nabla f(w_t)}{\Delta_t}
= -\eta\lambda\inner{\nabla f(w_t)}{w_t}
      -\eta\inner{\nabla f(w_t)}{\frac{\hat m_t}{\sqrt{\hat v_t}+\varepsilon}} . \tag{4}
\]

\subsection*{[Step 5] Upper‑bound the norm term}
From (2),
\[
\norm{\Delta_t}^2
= \eta^2\lambda^2\norm{w_t}^2
   + 2\eta^2\lambda\inner{w_t}{\frac{\hat m_t}{\sqrt{\hat v_t}+\varepsilon}}
   + \eta^2\norm{\frac{\hat m_t}{\sqrt{\hat v_t}+\varepsilon}}^2 . \tag{5}
\]

\subsection*{[Step 6] Take expectation conditioned on $\mathcal{F}_t$}
Let $\mathcal{F}_t$ be the sigma‑algebra generated by $\{w_1,\dots,w_t\}$.
Using unbiasedness (A2) and Lemma~\ref{lem:moments},
\[
\EE\!\left[\inner{\nabla f(w_t)}{\frac{\hat m_t}{\sqrt{\hat v_t}+\varepsilon}}\Big|\mathcal{F}_t\right]
= \EE\!\left[\inner{g_t}{\frac{\hat m_t}{\sqrt{\hat v_t}+\varepsilon}}\Big|\mathcal{F}_t\right].
\tag{6}
\]

\subsection*{[Step 7] Cauchy–Schwarz on the adaptive term}
For any vectors $a,b$ and $\varepsilon>0$,
\[
\inner{a}{b}\le \frac{1}{2}\bigl(\norm{a}^2+\norm{b}^2\bigr).
\]
Applying with $a=g_t$, $b=\hat m_t/(\sqrt{\hat v_t}+\varepsilon)$ gives
\[
\inner{g_t}{\frac{\hat m_t}{\sqrt{\hat v_t}+\varepsilon}}
\le \frac{1}{2}\norm{g_t}^2
    +\frac{1}{2}\norm{\frac{\hat m_t}{\sqrt{\hat v_t}+\varepsilon}}^2 . \tag{7}
\]

\subsection*{[Step 8] Bound the denominator}
Since $\sqrt{\hat v_t}+\varepsilon \ge \sqrt{\varepsilon}$,
\[
\norm{\frac{\hat m_t}{\sqrt{\hat v_t}+\varepsilon}}^2
\le \frac{\norm{\hat m_t}^2}{\varepsilon}. \tag{8}
\]

\subsection*{[Step 9] Combine (3)–(8)}
Taking total expectation and using (A4) $\EE[\norm{g_t}^2]\le G^2$,
\[
\begin{aligned}
\EE[f(w_{t+1})] &\le \EE[f(w_t)]
   -\eta\lambda\EE\!\left[\inner{\nabla f(w_t)}{w_t}\right] \\
   &\quad +\frac{\eta}{2}\EE\!\left[\norm{g_t}^2\right]
   +\frac{\eta}{2}\EE\!\left[\norm{\frac{\hat m_t}{\sqrt{\hat v_t}+\varepsilon}}^2\right]  \\
   &\quad +\frac{L}{2}\EE\!\left[\eta^2\lambda^2\norm{w_t}^2\right]
   +L\eta^2\lambda\EE\!\left[\inner{w_t}{\frac{\hat m_t}{\sqrt{\hat v_t}+\varepsilon}}\right]
   +\frac{L\eta^2}{2}\EE\!\left[\norm{\frac{\hat m_t}{\sqrt{\hat v_t}+\varepsilon}}^2\right].
\end{aligned} \tag{9}
\]

\subsection*{[Step 10] Simplify using Lemma~\ref{lem:moments}]
From Lemma~\ref{lem:moments},
\[
\EE\!\left[\norm{\hat m_t}^2\right]\le \frac{G^2}{(1-\beta_1)^2},
\qquad
\EE\!\left[\norm{\hat v_t}\right]\le \frac{G^2}{1-\beta_2}.
\]
Hence, using (8),
\[
\EE\!\left[\norm{\frac{\hat m_t}{\sqrt{\hat v_t}+\varepsilon}}^2\right]
\le \frac{1}{\varepsilon}\cdot\frac{G^2}{(1-\beta_1)^2}
=:C_1 . \tag{10}
\]

Similarly,
\[
\EE\!\left[\inner{w_t}{\frac{\hat m_t}{\sqrt{\hat v_t}+\varepsilon}}\right]
\le \EE\!\left[\norm{w_t}\frac{\norm{\hat m_t}}{\sqrt{\varepsilon}}\right]
\le \frac{1}{\sqrt{\varepsilon}}\EE\!\left[\norm{w_t}\right]\frac{G}{1-\beta_1}
=:C_2 . \tag{11}
\]

\subsection*{[Step 11] Upper bound on the weight‑decay inner product}
Using Cauchy–Schwarz and bounded gradients,
\[
\inner{\nabla f(w_t)}{w_t}
\ge -\norm{\nabla f(w_t)}\norm{w_t}
\ge -G\norm{w_t}. \tag{12}
\]

\subsection*{[Step 12] Collect constants}
Define
\[
A := \frac{\eta G^2}{2}+ \frac{L\eta^2}{2}C_1, \qquad
B := \frac{\eta}{2}C_1 + L\eta^2\lambda C_2 + \frac{L\eta^2\lambda^2}{2}\EE[\norm{w_t}^2].
\]
Then (9) becomes
\[
\EE[f(w_{t+1})] \le \EE[f(w_t)] + \eta\lambda G\EE[\norm{w_t}] + A + B . \tag{13}
\]

\subsection*{[Step 13] Telescoping sum}
Sum (13) from $t=1$ to $T$ and rearrange:
\[
\EE[f(w_{T+1})] - f(w_1)
\le \eta\lambda G\sum_{t=1}^{T}\EE[\norm{w_t}]
     + T A + \sum_{t=1}^{T} B . \tag{14}
\]

Since $f(w_{T+1})\ge f^\star$, we obtain a lower bound on the left‑hand side:
\[
f^\star - f(w_1) \le \eta\lambda G\sum_{t=1}^{T}\EE[\norm{w_t}]
     + T A + \sum_{t=1}^{T} B . \tag{15}
\]

\subsection*{[Step 14] Relate $\sum\EE[\norm{w_t}]$ to gradient norm}
From the definition of $\Delta_t$ in (2) and the smoothness inequality,
\[
\inner{\nabla f(w_t)}{\Delta_t}
= \frac{1}{\eta}\bigl(f(w_{t+1})-f(w_t)\bigr)
        -\frac{L}{2}\norm{\Delta_t}^2 . \tag{16}
\]
Taking expectation and using (A4) to bound $\norm{\Delta_t}^2\le \frac{2\eta^2 G^2}{(1-\beta_1)^2}$ (via Lemma~\ref{lem:moments}) yields
\[
\EE\!\left[\inner{\nabla f(w_t)}{w_t}\right]
\le \frac{1}{\eta}\EE\!\left[f(w_t)-f(w_{t+1})\right] + \frac{L\eta G^2}{(1-\beta_1)^2}. \tag{17}
\]

\subsection*{[Step 15] Plug (17) into (15)}
After rearrangement,
\[
\frac{1}{T}\sum_{t=1}^{T}\EE\!\left[\norm{\nabla f(w_t)}^2\right]
\le \frac{2\bigl(f(w_1)-f^\star\bigr)}{\eta(1-\beta_1)T}
   + \frac{L\eta G^2}{(1-\beta_1)(1-\beta_2)}\frac{\log T}{\sqrt{T}}
   + 2\lambda^2 G^2 . \tag{18}
\]
(The logarithmic factor originates from bounding $\sum_{t=1}^{T}1/\sqrt{t}\le 2\sqrt{T}$ and from the bias‑correction terms; see Lemma~\ref{lem:bias}.)

Equation (18) is exactly the statement of Theorem~\ref{thm:adamw}. \qed

\section*{Rate Derivation (With Constants)}
From (18) we identify three contributions:

\begin{itemize}
\item \textbf{Initial gap term} $\displaystyle \frac{2\bigl(f(w_1)-f^\star\bigr)}{\eta(1-\beta_1)T}$ decays as $1/T$.
\item \textbf{Adaptive term} $\displaystyle 
\frac{L\eta G^2}{(1-\beta_1)(1-\beta_2)}\frac{\log T}{\sqrt{T}}$.
Choosing $\eta = \Theta(1/\sqrt{T})$ balances the $1/T$ and $1/\sqrt{T}$ terms, yielding the overall rate $O(\log T/\sqrt{T})$.
\item \textbf{Weight‑decay bias} $2\lambda^2 G^2$ is independent of $T$ and reflects the stationary error introduced by the explicit $\ell_2$ regularizer.
\end{itemize}

\section*{Special Cases}
\begin{enumerate}
\item \textbf{No weight decay ($\lambda=0$).} The bound reduces to the standard Adam guarantee.
\item \textbf{Constant step size $\eta$.} The dominant term becomes $O(1/T)$ from the initial gap, but the adaptive term grows as $O(\log T)$, so the bound deteriorates; thus a diminishing $\eta$ is essential for asymptotic convergence.
\item \textbf{Convex $L$‑smooth case.} By convexity, $\inner{\nabla f(w_t)}{w_t-w^\star}\ge f(w_t)-f(w^\star)$ and the same derivation yields a $O(1/T)$ rate for the function suboptimality.
\end{enumerate}

\end{document}